% F1 fig:secant (Chapter 1, margin): the tangent at x and two secants, a diagram without data. The curve is
% exp(t) around x = 0 with a step of 0.7; any smooth curve would do, and the text does not name it. The forward
% secant joins the points at x and x + h; the central secant joins the points at x - h and x + h and is nearly
% parallel to the tangent: the one thing to notice, in accent. The lines end at staggered abscissas so that their
% labels do not collide.
\begin{tikzpicture}[x=17mm, y=10.5mm, font=\sffamily\scriptsize]
  \pgfmathsetmacro{\stepH}{0.7}
  \pgfmathsetmacro{\fMinus}{exp(-\stepH)}
  \pgfmathsetmacro{\fPlus}{exp(\stepH)}
  \pgfmathsetmacro{\slopeF}{(\fPlus - 1) / \stepH}              % slope of the forward secant
  \pgfmathsetmacro{\slopeC}{(\fPlus - \fMinus) / (2 * \stepH)}  % slope of the central secant
  % the axis with the three nodes
  \draw[thin, black!55] (-1.2, -0.3) -- (1.25, -0.3);
  \foreach \nodeX/\nodeLabel in {-\stepH/{$x - h$}, 0/{$x$}, \stepH/{$x + h$}} {
    \draw[thin, black!55] (\nodeX, -0.3) -- ++(0, -0.05)
      node[below, black!70, inner sep=1.5pt] {\nodeLabel};
    \draw[densely dotted, black!40] (\nodeX, -0.3) -- (\nodeX, {exp(\nodeX)});
  }
  % the function
  \draw[line width=0.8pt, black!80, domain=-1.2:1.22, samples=60, smooth] plot (\x, {exp(\x)})
    node[left, black!80, inner sep=2pt] {$f$};
  % the tangent at x: slope f'(x) = 1
  \draw[line width=0.6pt, black!80] (-1.15, {1 - 1.15}) -- (0.95, {1 + 0.95})
    node[right, inner sep=1.5pt] {tangent};
  % the central secant through x - h and x + h
  \draw[line width=0.6pt, accent] (-1.15, {\fMinus + (-1.15 + \stepH) * \slopeC})
    -- (1.05, {\fMinus + (1.05 + \stepH) * \slopeC})
    node[right, inner sep=1.5pt] {central};
  % the forward secant through x and x + h
  \draw[line width=0.6pt, muted] (-0.3, {1 - 0.3 * \slopeF}) -- (1.15, {1 + 1.15 * \slopeF})
    node[right, inner sep=1.5pt] {forward};
  % the sampled points
  \fill[muted] (\stepH, \fPlus) circle (1.1pt);
  \fill[accent] (-\stepH, \fMinus) circle (1.1pt);
  \fill[black!80] (0, 1) circle (1.1pt);
\end{tikzpicture}
